\NormalFont\ShortTitle{1 Pedro}
{\MT Primeira Carta Universal de Pedro

\par }\ChapOne{1}{\PP \VerseOne{1}Pedro, apóstolo de Jesus Cristo, aos forasteiros que estão espalhados em Ponto, Galácia, Capadócia, Ásia e Bitínia,
\VS{2}escolhidos segundo o pré-conhecimento de Deus Pai, por meio da santificação do Espírito, para a obediência e a aspersão do sangue de Jesus Cristo. Que a graça e a paz vos sejam multiplicadas.
\VS{3}Bendito seja o Deus e Pai de nosso Senhor Jesus Cristo. Segundo sua grande misericórdia, ele nos regenerou para uma esperança viva, por meio da ressurreição de Jesus Cristo dentre os mortos.
\VS{4}E o resultado disso é uma herança incorruptível, incontaminável, e que não pode ser enfraquecida. Ela está guardada nos céus para vós,
\VS{5}que pela fé estais guardados no poder de Deus para a salvação, pronta para se revelar no último tempo.
\VS{6}Nisso vos alegrais, ainda que agora seja necessário que estejais por um pouco
{\ADD{de tempo}} entristecidos com várias provações,
\VS{7}para que a prova de vossa fé, muito mais preciosa que o ouro que perece (mesmo sendo provado pelo fogo), seja achada em louvor, honra, e glória, na revelação de Jesus Cristo.
\VS{8}A ele, sem terdes visto, vós o amais. Ainda que não estejais o vendo agora, mas crendo nele, vós estais alegres com júbilo indescritível e glorioso.
\VS{9}Assim estais alcançando o resultado da vossa fé: a salvação das
{\ADD{vossas}} almas.
\VS{10}Dessa salvação os profetas que profetizaram a respeito da graça
{\ADD{direcionada}} a vós procuraram entender, e com empenho investigaram,
\VS{11}procurando saber qual era o tempo
\FTNT{*}{{\FR 1:11 }
o tempo
ou: a pessoa
} ou ocasião que o Espírito de Cristo dentro deles indicava, quando dava prévio testemunho dos sofrimentos de Cristo e das glórias seguintes.
\VS{12}A eles foi revelado que não foi para eles mesmos, mas sim para nós, que eles prestaram serviço com estas coisas. E agora, elas vos foram anunciadas pelos que, pelo Espírito Santo enviado do céu, vos pregaram o Evangelho. Tais coisas os anjos desejam atentamente observar.
\VS{13}Portanto, preparai as vossas mentes, sede sóbrios, e esperai completamente na graça que vos será trazida na revelação de Jesus Cristo.
\FTNT{†}{{\FR 1:13 }
preparai vossas mentes
lit. revestindo os lombos de vossa mente
}
\VS{14}E, como filhos obedientes, não vos conformeis com os maus desejos do passado, de quando éreis ignorantes.
\FTNT{‡}{{\FR 1:14 }
de quando éreis ignorantes
lit. em vossa ignorância
}
\VS{15}Ao contrário; assim como aquele que vos chamou é santo, sede vós também santos em todo o
{\ADD{vosso}} comportamento.
\VS{16}Pois está escrito: Sede santos, porque eu sou santo.
{\RQ{Levítico 11:44-45
}}
\VS{17}E, se vós chamais de Pai aquele que, sem acepção de pessoas, julga segundo a obra de cada um, vivei em temor o tempo de vossa peregrinação;
\VS{18}e sabei que não foi por coisas destrutíveis,
{\ADD{como}} prata ou ouro, que fostes resgatados da vossa vã maneira de viver que recebestes dos
{\ADD{vossos}} pais,
\VS{19}mas sim, pelo sangue precioso, como de um cordeiro sem falha e sem contaminação: o
{\ADD{sangue}} de Cristo.
\VS{20}Ele já era conhecido
\FTNT{§}{{\FR 1:20 }
conhecido
isto é, conhecido o seu papel como o Salvador. Trad. alt: predestinado
} desde antes da fundação do mundo, mas foi manifesto nestes últimos tempos por causa de vós,
\VS{21}que, por meio dele, credes em Deus, que o ressuscitou dos mortos e lhe deu glória, para que a vossa fé e esperança estejam em Deus.
\VS{22}Uma vez que purificastes as vossas almas na obediência da verdade por meio do Espírito, para o sincero
\FTNT{*}{{\FR 1:22 }
sincero
ou: não-hipócrita
} amor fraternal, amai-vos intensamente uns aos outros com um coração puro,
\VS{23}visto que nascestes de novo, não de semente que perece, mas sim imperecível, pela viva palavra de Deus, que permanece para sempre.
\VS{24}Pois toda carne
\FTNT{†}{{\FR 1:24 }
toda carne
i.e. todos os seres humanos
} é como a erva, e toda glória humana como a flor da erva. A erva se seca, e sua flor cai;
\VS{25}mas a palavra do Senhor permanece para sempre.E esta é a palavra que foi evangelizada entre vós.
{\RQ{Isaías 40:6,8
}}

\par }\Chap{2}{\PP \VerseOne{1}Abandonai, portanto, toda malícia, toda enganação, fingimentos, invejas, e todas as falas maldosas,
\VS{2}e desejai ansiosamente, como bebês recém-nascidos, o leite da Palavra, não falsificado, para que por meio dele estejais crescendo;
\VS{3}se, de fato, já experimentastes que o Senhor é bom.
\VS{4}Aproximai-vos dele,
{\ADD{que é}} uma pedra viva, rejeitada pelos seres humanos, mas escolhida e preciosa para Deus.
\VS{5}Também vós, como pedras vivas, sois edificados como casa espiritual
{\ADD{e}} sacerdócio santo, para oferecerdes sacrifícios espirituais, agradáveis a Deus por meio de Jesus Cristo.
\VS{6}Por isso também está contido na Escritura: Eis que eu ponho em Sião uma pedra principal de esquina, escolhida
{\ADD{e}} preciosa. Quem nela crer de maneira nenhuma será envergonhado.
\FTNT{*}{{\FR 2:6 }
ou: frustrado
}
{\RQ{Isaías 28:16
}}
\VS{7}Assim, para vós, que credes, ela é preciosa. Mas para os rebeldes: A pedra que os construtores rejeitaram, essa se tornou a principal de esquina,
{\RQ{Salmo 118:22
}}
\VS{8}assim como pedra de tropeço, e pedra que causa queda. Eles tropeçam na palavra por serem rebeldes. Para isso eles foram destinados.
{\RQ{Isaías 8:14
}}
\VS{9}Mas vós sois a geração escolhida, o sacerdócio real, a nação santa, o povo adquirido; a fim de que anuncieis as virtudes daquele que vos chamou das trevas para a sua maravilhosa luz.
\VS{10}Antes, vós não éreis povo, mas agora sois povo de Deus.
{\ADD{Antes}} , não havíeis recebido misericórdia, mas agora recebestes misericórdia.
\VS{11}Amados, como a peregrinos e estrangeiros, eu vos peço que vos abstenhais dos desejos carnais, que batalham contra a alma,
\VS{12}e que tenhais um bom comportamento vosso entre os pagãos; para que, naquilo que de vós, como de malfeitores, falam mal, no dia da visitação glorifiquem a Deus por causa das
{\ADD{vossas}} boas obras que tiverem visto.
\VS{13}Portanto sujeitai-vos a toda autoridade humana, por causa do Senhor; seja ao rei, como superior;
\VS{14}seja aos governantes, como enviados por ele, com o objetivo de castigar os malfeitores, e de conceder honra aos que fazem o bem.
\VS{15}Pois esta é a vontade de Deus, que, ao fazerdes o bem, caleis a ignorância dos homens tolos.
\VS{16}{\ADD{Comportai-vos}} como pessoas livres, mas não useis a liberdade como pretexto para a malícia. Em vez disso,
{\ADD{sede}} como servos de Deus.
\VS{17}Honrai a todos: amai aos irmãos, temei a Deus, honrai ao rei.
\VS{18}Servos, sujeitai-vos com todo temor aos
{\ADD{vossos}} senhores, não somente aos bons e brandos, mas também aos que maltratam.
\VS{19}Pois coisa agradável é se alguém, por causa da consciência a respeito de Deus, experimente dores, sofrendo injustamente.
\VS{20}Afinal, que mérito há em suportar serdes espancados por cometerdes pecado? Contudo, se, quando fazeis o bem, sois afligidos e suportais, isso é a Deus.
\VS{21}Pois para isto fostes chamados, porque também Cristo sofreu por nós, deixando-nos exemplo, para que sigais os seus passos.
\VS{22}Ele não cometeu pecado, nem engano foi achado em sua boca.
\VS{23}Quando o insultavam, ele não insultava de volta. Quando sofria, ele não ameaçava; em vez disso, entregava-se ao que julga de maneira justa.
\VS{24}Ele levou nossos pecados em seu próprio corpo sobre o madeiro; para que nós, estando mortos para os pecados, vivamos para a justiça. Pela ferida dele fostes sarados.
\VS{25}Porque vós éreis como ovelhas desviadas do caminho; mas agora vós estais convertidos ao Pastor e Supervisor de vossas almas.

\par }\Chap{3}{\PP \VerseOne{1}Semelhantemente vós, esposas, estejais sujeitas aos próprios maridos, para que, ainda que alguns não obedeçam à Palavra, por meio do comportamento das esposas, sem palavra, sejam ganhos,
\VS{2}quando eles observarem o vosso comportamento puro e com temor.
\VS{3}A vossa beleza não seja a exterior,
{\ADD{como}} tranças de cabelos, joias de ouro, ou vestuário,
\VS{4}mas sim, a pessoa interior no coração, o
{\ADD{embelezamento}} incorruptível de um espírito manso e tranquilo, que é precioso diante de Deus.
\VS{5}Pois assim também se embelezavam antigamente as mulheres santas, que esperavam em Deus, e eram sujeitas aos seus maridos.
\VS{6}Dessa maneira Sara obedecia a Abraão, chamando-lhe de senhor. Vós sois filhas dela se fizerdes o bem e não temerdes espanto algum.
\VS{7}Semelhantemente vós maridos, convivei
{\ADD{com elas}} com entendimento, dando honra à mulher, como a vaso mais fraco, como sendo
{\ADD{elas e vós}} juntamente herdeiros da graça da vida; a fim de que as vossas orações não sejam impedidas.
\VS{8}E por fim, sede todos de uma mesma mentalidade, compassivos, tenhais amor pelos irmãos, e sede misericordiosos e benevolentes.
\VS{9}Não retribuais o mal com mal, ou insulto com insulto; ao contrário, bendizei, sabendo que para isto fostes chamados, para que herdeis a bênção.
\VS{10}Pois: quem quer amar a vida e ver dias bons, refreie sua língua do mal, e que seus lábios não falem engano.
\VS{11}Afaste-se do mal, e faça o bem; busque a paz, e siga-a.
\VS{12}Porque os olhos do Senhor estão sobre os justos, e seus ouvidos
{\ADD{atentos}} às suas orações; mas a face do Senhor é contra os que fazem o mal.
{\RQ{Salmos 34:12-16
}}
\VS{13}E quem vos fará mal se fordes adeptos do bem?
\VS{14}Porém, mesmo se sofrerdes por causa da justiça, benditos
\FTNT{*}{{\FR 3:14 }
benditos
ou: bem-aventurados
} sois. Não tenhais medo deles, nem vos perturbeis.
\VS{15}Mas santificai a Deus como Senhor em vossos corações; e estai sempre preparados para responder com mansidão e respeito a todo aquele que vos pedir a razão da esperança que há em vós,
\VS{16}tendo uma boa consciência; para que, naquilo que falam mal de vós como malfeitores, os que insultam o vosso bom comportamento em Cristo sejam envergonhados.
\VS{17}Pois é melhor que sofrais fazendo o bem, se
{\ADD{assim}} a vontade de Deus quer, do que fazendo o mal.
\VS{18}Porque também Cristo sofreu uma vez pelos pecados, o justo pelos injustos; para que nos levasse a Deus.
{\ADD{Ele estava}} , de fato, morto na carne, mas vivificado pelo Espírito,
\VS{19}no qual ele também foi pregar aos espíritos em prisão.
\VS{20}{\ADD{Estes são}} os que antigamente foram rebeldes, quando uma vez a paciência de Deus aguardava nos dias de Noé, enquanto era preparada a arca. Nela, poucas almas (isto é, oito) foram salvas por meio da água.
\VS{21}Esta é uma representação do batismo, que agora também nos salva, não como remoção da sujeira do corpo, mas sim como o pedido
\FTNT{†}{{\FR 3:21 }
pedido
ou: compromisso. A palavra usada no texto original grego tem significado obscuro para o contexto da frase
} de boa consciência a Deus, por meio da ressurreição de Jesus Cristo.
\VS{22}Ele está à direita de Deus, depois de haver subido ao céu; e os anjos, as autoridades, e os poderes estão sob o seu comando.

\par }\Chap{4}{\PP \VerseOne{1}Visto que Cristo sofreu por nós na carne, armai-vos também com o mesmo modo de pensar, pois quem sofreu na carne terminou com o pecado;
\VS{2}para que, no resto do tempo na carne, não viva mais segundo os maus desejos humanos, mas sim segundo a vontade de Deus.
\VS{3}Pois já basta, no tempo passado da vida, termos feito a vontade dos pagãos,
\FTNT{*}{{\FR 4:3 }
pagãos
ou: gentios
} e andado em promiscuidades, desejos maliciosos, bebedeiras, orgias, festas de embriaguez, e abomináveis idolatrias.
\VS{4}Como eles estranham vós não agirdes como eles, que correm em busca da extrema devassidão, por isso falam mal de vós.
\FTNT{†}{{\FR 4:4 }
agirdes … devassidão
lit. correrdes [com eles] à mesma extrema devassidão
}
\VS{5}Eles terão de prestar contas ao que está pronto para julgar os vivos e os mortos.
\VS{6}Pois por isso o evangelho também foi pregado aos mortos,
\FTNT{‡}{{\FR 4:6 }
i.e. Os que agora estão mortos, mas que receberam o evangelho enquanto ainda estavam vivos
} para que, mesmo que tenham sido julgados na carne segundo os critérios humanos,
\FTNT{§}{{\FR 4:6 }
segundo os critérios humanos
lit. segundo os seres humanos
} vivam em espírito segundo os critérios de Deus.
\FTNT{*}{{\FR 4:6 }
segundo os critérios de Deus
lit. segundo Deus. Outra alternativa seria: como Deus vive [em Espírito]
}
\VS{7}O fim de todas as coisas está próximo; portanto sede sóbrios, e vigiai em orações.
\VS{8}Mas, acima de tudo, tende fervoroso amor uns para com os outros; porque o amor cobrirá uma multidão de pecados.
\VS{9}Proporcionai hospitalidade uns aos outros, sem murmurações.
\VS{10}Cada um sirva aos outros segundo o dom que recebeu, como bons administradores da variada graça de Deus.
\VS{11}Se alguém fala, que seja como as palavras de Deus; se alguém serve,
{\ADD{que sirva}} segundo a força que Deus dá; a fim de que em tudo Deus seja glorificado por meio de Jesus Cristo. E a ele pertencem a glória e o poder para todo o sempre, Amém!
\VS{12}Amados, não estranheis o fogo ardente que vem sobre vós para vos provar, como se algo estranho estivesse vos acontecendo.
\VS{13}Em vez disso, alegrai-vos em serdes participantes das aflições de Cristo; para que, também na revelação de sua glória vos alegreis com júbilo.
\VS{14}Se vós sois insultados por causa do nome de Cristo, benditos
\FTNT{†}{{\FR 4:14 }
ou: bem-aventurados
} sois; porque o Espírito da glória e de Deus repousa sobre vós; o qual por eles é insultado, mas por vós é glorificado.
\VS{15}Porém nenhum de vós sofra como homicida, ladrão, malfeitor, ou como alguém que se intromete em assuntos alheios;
\FTNT{‡}{{\FR 4:15 }
alguém que se intromete em assuntos alheios
a palavra no original tem significado obscuro. Ou: defraudador
}
\VS{16}Mas se
{\ADD{sofre}} como cristão, não se envergonhe; ao contrário, glorifique a Deus por isso;
\VS{17}Pois já é tempo do julgamento começar pela casa de Deus; e se é primeiro conosco, qual será o fim daqueles que são desobedientes ao Evangelho de Deus?
\VS{18}E se o justo se salva com dificuldade, em que situação aparecerá o ímpio e pecador?
{\RQ{Provérbios 11:31
}}
\VS{19}Portanto, os que sofrem segundo a vontade de Deus, confiem suas almas ao fiel Criador, fazendo o bem.

\par }\Chap{5}{\PP \VerseOne{1}Aos anciãos
\FTNT{*}{{\FR 5:1 }
ou: presbíteros
} da igreja que estão entre vós, eu, que sou ancião como eles, e testemunha dos sofrimentos de Cristo, e participante da glória que se revelará, seriamente peço:
\VS{2}Pastoreai o rebanho de Deus que está entre vós, não por obrigação, mas sim voluntariamente; não por ganância, mas sim com sincera prontidão;
\VS{3}nem como que dominando os que estão aos vossos cuidados,
\FTNT{†}{{\FR 5:3 }
os que estão aos vossos cuidados
lit. os reservados, ou separados. No vocabulário bíblico, é frequentemente usado para indicar propriedades e possessões, em particular propriedades de terra. Por esse motivo, traduções mais antigas põem: propriedades [de Deus] (ou algo semelhante). Outra interpretação é a de que se refere aos anciãos, referidos por Pedro como “vós”, e de que a reserva ou separação não é para a obtenção ou posse, mas tão somente para a responsabilidade e o cuidado. Isto é, refere-se aos cristãos que os anciãos devem cuidar. Esta última interpretação é a usada na tradução.
} mas sim como exemplos ao rebanho.
\VS{4}E quando o Pastor Principal aparecer, vós recebereis a indestrutível coroa da glória.
\VS{5}Semelhantemente vós, jovens, sede sujeitos aos anciãos; e todos vós sede sujeitos uns aos outros, revestidos de humildade; pois: Deus resiste aos soberbos, mas dá graça aos humildes.
{\RQ{Provérbios 3:34
}}
\VS{6}Humilhai-vos, pois, debaixo da poderosa mão de Deus, para que ele vos exalte no tempo adequado,
\VS{7}lançando sobre ele toda a vossa ansiedade; porque ele tem cuidado de vós.
\VS{8}Sede sóbrios! Vigiai! Pois o vosso adversário, o diabo, anda ao redor, rugindo como um leão, buscando a quem possa devorar.
\VS{9}Resisti a ele, firmes na fé; sabendo que as mesmas aflições acontecem com os vossos irmãos
\FTNT{‡}{{\FR 5:9 }
os vossos irmãos
lit. a vossa irmandade
} no mundo.
\VS{10}E o Deus de toda graça, que nos chamou para a sua eterna glória em Cristo Jesus, depois de sofrerdes um pouco, ele vos aperfeiçoe, confirme, fortifique, e estabeleça fundamento.
\VS{11}A ele sejam a glória e o poder para todo o sempre. Amém!
\VS{12}Por Silvano, vosso fiel irmão, como o considero, eu vos escrevi brevemente, exortando e dando testemunho de que esta é a verdadeira graça de Deus, na qual estais.
\VS{13}Saúda-vos a
{\ADD{igreja}} escolhida convosco, que está na Babilônia, e também o meu filho Marcos.
\VS{14}Saudai-vos uns aos outros com beijo de amor. A paz seja com todos vós que estais em Cristo Jesus. Amém.\par }